% !TeX spellcheck = ru_RU
% !TEX root = vkr.tex

\section*{Заключение}

В рамках учебной практики второго семестра получены следующие результаты:
\begin{itemize}
    \item выполнен обзор статического кодирования Хаффмана и его применения в
          составе существующих алгоритмов сжатия;
    \item разработан консольный архиватор на языке~C с документированным
          бинарным форматом и поддержкой текстовых и бинарных файлов;
    \item корректность реализации проверена модульными и интеграционными
          тестами: в воспроизведённом запуске успешно прошли 6 из 6 модульных
          наборов тестов и обработаны 6 из 6 интеграционных входных файлов;
    \item проведены измерения и приведены результаты для синтетических и
          типовых файлов; подтверждено эффективное сжатие избыточных данных,
          отсутствие выигрыша для равномерно случайного входа и зависимость
          результата от внутренней структуры реального файла.
\end{itemize}

Исходный код, документация по сборке, описание формата и сценарии экспериментов
опубликованы в репозитории проекта:
\url{https://github.com/RomanKazz/huffman-archiver}.

Для проекта настроены автоматические проверки сборки, тестов, форматирования и
статического анализа при изменениях в репозитории. Дальнейшим развитием работы
может быть фиксация входных данных эксперимента для полностью детерминированного
сравнения, сопоставление результата с форматами, использующими \Deflate{}, и
расширение набора измерений более крупными пользовательскими файлами разных
форматов.
